\section{Introdução}
\label{sec:introducao}

A crescente penetração da geração fotovoltaica nas redes elétricas tem sido
acompanhada por desafios relacionados à intermitência e à variabilidade
dessa fonte, que comprometem a estabilidade do fornecimento de energia e a
qualidade de tensão em sistemas de pequeno e médio porte
\citep{ayaz2014improved}. Para mitigar esses efeitos, sistemas de
armazenamento de energia em baterias (\textit{Battery Energy Storage
System} -- BESS) têm sido amplamente incorporados a instalações
fotovoltaicas, atuando como reguladores de potência e garantindo o
suprimento contínuo da carga mesmo sob variações bruscas de irradiância.

A integração entre o painel fotovoltaico, o banco de baterias e a carga é
tradicionalmente realizada por meio de conversores CC-CC independentes, um
para cada fonte, conectados a um barramento comum. Essa abordagem, embora
simples, implica maior número de componentes, perdas adicionais de
conversão e custo elevado. Como alternativa, conversores CC-CC de múltiplas
portas têm sido propostos para integrar essas três interfaces em um único
estágio de potência, reduzindo a complexidade do sistema sem comprometer a
flexibilidade no gerenciamento do fluxo de energia \citep{zhang2025adaptive}.

Por outro lado, a operação desses conversores envolve o acoplamento entre
múltiplas portas, incertezas paramétricas e perturbações de carga não
lineares, fatores que reduzem a eficácia de controladores clássicos de
ganho fixo, como o PID, frente a variações bruscas de operação
\citep{zhang2025adaptive}. Técnicas de controle preditivo baseado em modelo
têm-se mostrado uma alternativa robusta para esse cenário, por incorporar
diretamente, na própria formulação do problema de controle, as restrições
físicas do sistema e suas múltiplas variáveis. Dentro dessa classe de
controladores, o Controle Preditivo Baseado em Modelo com Conjunto Finito
de Controle (\textit{Finite Control Set Model Predictive Control} --
FCS-MPC) explora diretamente a natureza discreta dos conversores
eletrônicos de potência, prescindindo da etapa de modulação por largura de
pulso (PWM) e permitindo respostas dinâmicas rápidas mesmo em sistemas
multivariáveis \citep{ferreira2018finite, wang2023fcsmpc}.

Neste trabalho, adota-se a topologia de conversor de três portas do tipo
VR-BESS (\textit{Voltage Regulator -- Battery Energy Storage System}), que
integra a porta fotovoltaica, a porta de armazenamento e a porta de carga
utilizando apenas duas chaves semicondutoras e um indutor compartilhado,
operando como conversor \textit{Buck} durante o carregamento da bateria e
como \textit{Boost} na regulação da tensão de saída. O número reduzido de
chaves resulta em um conjunto de estados de comutação naturalmente pequeno,
o que torna essa topologia particularmente adequada à aplicação direta do
FCS-MPC, cujo custo computacional cresce com o número de estados avaliados
a cada período de amostragem.

O objetivo deste trabalho é desenvolver e simular o controle do conversor
VR-BESS por meio da técnica FCS-MPC, a partir da modelagem em espaço de
estados do conversor e do painel fotovoltaico associado. Além disso, serão
estudados observadores de estado como alternativa para reduzir a
quantidade de sensores físicos necessários à implementação do
controlador, contribuindo para uma solução de controle mais simples e de
menor custo de instrumentação. Como etapa prévia à síntese do controlador,
este trabalho apresenta a modelagem e a validação em malha aberta do
conversor VR-BESS e de seus subsistemas de painel fotovoltaico e banco de
baterias, implementados como bibliotecas próprias no Simscape Electrical.
Essa validação estabelece a base de simulação sobre a qual o controlador
FCS-MPC será desenvolvido e avaliado, em cenários de variação de
irradiância, carga e referência de tensão, na etapa subsequente do
trabalho.
